\chapter{Introducción}

\section{Contexto y motivación}

Los equipos de ingeniería toman decisiones continuamente, tales como qué arquitectura adoptar, qué
riesgos aceptar, qué tareas asignar y qué preguntas quedan abiertas. Una parte importante de ese
proceso ocurre en reuniones técnicas, donde la conversación directa facilita la deliberación y la
toma de acuerdos. Sin embargo, el conocimiento generado en esas reuniones rara vez queda
registrado de manera estructurada: termina disperso en notas informales, hilos de mensajería
instantánea y en la memoria de los participantes, sin trazabilidad del razonamiento que condujo
a cada decisión.

Esta brecha entre el conocimiento producido y el conocimiento preservado tiene consecuencias
prácticas: cuando un nuevo miembro del equipo necesita entender por qué se eligió una base de
datos en particular, o cuando un riesgo identificado meses atrás resurge en un contexto diferente,
no existe un registro consultable que responda a esas preguntas. La ausencia de un grafo de
decisiones explícito obliga a reconstruir el contexto a partir de fuentes fragmentadas, con el
coste de tiempo y el riesgo de inconsistencia que ello conlleva.

Existen herramientas que transcriben y resumen reuniones automáticamente, pero su salida es texto
no estructurado: un resumen o una transcripción, no un conjunto de entidades consultables. Ninguna
de ellas distingue entre una decisión, un riesgo o una tarea; ninguna rastrea cómo una decisión
tomada en marzo contradice o matiza otra tomada en junio. La estructuración del conocimiento, su
tipado y la trazabilidad de su evolución a lo largo del tiempo quedan fuera de su alcance.

Los \ac{LLM} han transformado las capacidades de procesamiento del lenguaje natural. Desde la
introducción de la arquitectura Transformer~\cite{vaswani2017attention}, los modelos de lenguaje
han escalado en capacidad hasta alcanzar rendimiento notable en tareas de comprensión y generación
de texto~\cite{brown2020gpt3}. Esta evolución ha abierto la puerta a sistemas capaces de extraer
información estructurada a partir de texto no estructurado de forma fiable, sin necesidad de
entrenamiento supervisado específico por dominio.

Al mismo tiempo, la técnica de \ac{RAG}~\cite{lewis2020rag} ha demostrado que combinar la
búsqueda en bases de conocimiento externas con la generación de modelos de lenguaje mejora
significativamente la precisión y trazabilidad de las respuestas. Esta combinación resulta
especialmente relevante en contextos donde el conocimiento evoluciona a lo largo del tiempo y
la coherencia entre distintas fuentes es crítica.

El presente \ac{TFM} presenta \textbf{Seshat}: un sistema que ingiere grabaciones de reuniones
técnicas, extrae entidades estructuradas mediante una cadena de agentes LLM, y las persiste en
una base de conocimiento con forma de grafo. El diseño incorpora un paso de revisión humana,
asistida por puntuación de confianza; un arnés de evaluación reproducible que permite comparar
el rendimiento del sistema cuando se introducen cambios, y una capa de observabilidad basada en
MLflow. El resultado es un sistema completo, desplegable mediante Docker Compose, orientado a
equipos de ingeniería que necesitan documentar y consultar el conocimiento técnico generado en
sus reuniones.

\section{Objetivos}

\subsection{Objetivo general}

El objetive principal del proyecto es diseñar e implementar un sistema de extracción automática
de conocimiento a partir de grabaciones de reuniones técnicas, capaz de identificar elementos
clave de forma estructurada, gestionar las relaciones entre ellas a lo largo del tiempo mediante
un grafo de conocimiento, e incorporar revisión humana asistida y evaluación calibrada como
requisitos de calidad.

\subsection{Objetivos específicos}

Por simplicidad, se descompone el objetivo general en los siguientes objetivos específicos:

\begin{enumerate}
  \item Diseñar un flujo multi-agente que extraiga entidades estructuradas de cuatro
        tipos a partir de transcripciones de reuniones e infiera las relaciones semánticas entre
        ellas y el conocimiento preexistente.

  \item Incorporar revisión humana asistida por puntuación de confianza, de modo que el revisor
        pueda concentrar su atención en los nodos de menor certeza y el sistema pueda operar en
        modo automático cuando la confianza lo permita.

  \item Construir un marco de evaluación con corpus etiquetado que mida el rendimiento mediante
        de forma reproducible, usando métricas como \textit{precision}, \textit{recall} o derivadas
        de ellas y bloquee el procesamiento de datos reales si no se alcanza la calidad definida.

  \item Desarrollar una \ac{API} REST y una interfaz web que cubran el ciclo de vida completo del
        trabajo (envío, transcripción, extracción, revisión y consulta del grafo) sin requerir
        acceso directo a la capa de persistencia.

  \item Desplegar el sistema como servicios contenerizados con observabilidad integrada, de modo
        que el rendimiento y el coste sean medibles y comparables entre versiones.
\end{enumerate}

\section{Alcance de la solución}
\label{sec:alcance}

El sistema admite dos tipos de entrada: archivos de audio (\texttt{.mp3}, \texttt{.wav},
\texttt{.m4a}) transcritos mediante un servicio externo, y transcripciones preformateadas
(\texttt{.yaml} o \texttt{.json}). Ambas convergen en una representación interna común antes
de entrar en el proceso de extracción.

La extracción cubre cuatro tipos de nodo y siete tipos de relación semántica entre ellos,
detallados en el capítulo~\ref{ch:requisitos-diseno}. La base de conocimiento sigue un modelo
de escritura por estado: el contenido de un nodo aprobado es estable por diseño, aunque
operadores y administradores pueden corregirlo mediante un mecanismo de anulación explícita
que registra el motivo y el autor del cambio. El único campo que cambia de forma automática
es el estado del nodo, que en esta primera versión se actualiza cuando una reunión posterior
establece una relación de supersesión o enmienda.

El sistema incluye una interfaz gráfica para la revisión humana de los nodos extraídos, un
mecanismo de aprobación masiva por umbral de confianza y la posibilidad de que el operador
edite nodos manualmente durante la revisión. Las relaciones inferidas se escriben
automáticamente tras la aprobación de los nodos. La observabilidad está cubierta mediante
\href{https://mlflow.org}{MLflow}~\cite{zaharia2018mlflow}: cada ejecución registra el uso
de tokens, la latencia por etapa y los resultados de las métricas de evaluación.

El sistema distingue tres tipos de trabajo: el \textit{pipeline} de ingesta y extracción,
que constituye el flujo principal descrito en este documento; los trabajos de
inicialización\footnote{%
  Los trabajos de inicialización, \textit{init}, poblarían la base de conocimiento a partir
  de una fuente distinta de una reunión —por ejemplo, un corpus documental existente— antes
  de la puesta en marcha del sistema. Su implementación se aplaza como trabajo futuro
  (capítulo~\ref{ch:conclusiones}, sección~\ref{sec:trabajo-futuro}).%
}, que preparan la base de conocimiento antes de que el sistema entre en operación; y las
operaciones manuales del operador, como la adición directa de nodos durante la revisión. En
lo que sigue, el término \textit{pipeline} se reserva para el flujo de ingesta y extracción.

\bigskip
\noindent Esta versión inicial (\acs{MVP}) deja fuera de su alcance un conjunto de capacidades
que se agrupan en cuatro categorías según la razón de su aplazamiento. La primera es el
\textbf{endurecimiento para producción}: el despliegue en infraestructura cloud, con su definición
mediante \ac{IaC}, y la autenticación con un proveedor de identidad externo
(\textit{SSO}) exceden lo que un prototipo de investigación necesita demostrar. La segunda son las
\textbf{extensiones del núcleo generativo}, funcionalidades afines a los objetivos pero pospuestas
por tiempo o por depender de datos de producción: la detección y anonimización de información
personal identificable (PII), la recuperación agéntica como alternativa al \textit{pipeline}
\ac{RAG}~\cite{yao2023react}, o una señal de confianza complementaria basada en inferencia de
lenguaje natural\footnote{%
  La \textit{inferencia de lenguaje natural} (\ac{NLI}) es la tarea de determinar si una hipótesis
  se puede deducir (\textit{entailment}), contradice (\textit{contradiction}) o es neutral respecto
  a una premisa dada. Aplicada a la puntuación de confianza, un modelo NLI verifica si la descripción
  del nodo extraído está respaldada por su cita fuente, aún usando sinónimos. Es, por ahora, una
  propuesta analizada pero no implementada: no se ha verificado si aportaría una señal realmente
  independiente de las heurísticas ya existentes.%
} (\acs{NLI}). La tercera es la \textbf{profundización del grafo de conocimiento}, trabajo de
modelado de datos más que de \ac{IAG}: la revisión humana de las relaciones inferidas o el uso de un
\textit{backend} de grafo nativo como \href{https://neo4j.com}{Neo4j}. La cuarta abarca la
\textbf{integración y nuevas fuentes de entrada}: la entrada de vídeo, la
diarización\footnote{%
  La \textit{diarización} (\textit{speaker diarization}) es el proceso de segmentar una grabación
  de audio en intervalos etiquetados por locutor, respondiendo a la pregunta «¿quién habló
  cuándo?».%
} de locutores, y la inicialización de la base desde un corpus documental existente —por ejemplo
desde \href{https://www.notion.so}{Notion} o \href{https://www.atlassian.com/software/confluence}{Confluence}—.
El capítulo~\ref{ch:conclusiones} retoma y detalla estas líneas como trabajo futuro
(sección~\ref{sec:trabajo-futuro}).
